\chapter{Presenting Quantitative Research}
\label{ch:presenting}

\begin{projectconnection}
This is preparation for Week~20---the final presentation.  Everything in this
chapter maps directly to Phase~5 deliverables: the research report (Task~23)
and the presentation itself (Task~24).  The framing advice here draws on how
the best academic papers and desk presentations communicate quantitative
results.  You have built infrastructure
(Chapter~\ref{ch:validation}), engineered features
(Chapter~\ref{ch:features}), labeled targets
(Chapter~\ref{ch:labeling}), validated honestly
(Chapter~\ref{ch:validation}), and backtested with costs
(Chapter~\ref{ch:backtesting}).  This chapter teaches you how to
\emph{communicate} all of that so a desk head trusts your work in 20 minutes.
\end{projectconnection}

% ══════════════════════════════════════════════════════════════
\section{Framing: Causal Hypothesis + ML Testing}
\label{sec:presenting:framing}
% ══════════════════════════════════════════════════════════════

\begin{prereq}[Why Framing Determines Reception]
A quantitative result is not evaluated in a vacuum.  The audience---whether a
trading desk, a research seminar, or a hiring committee---decides within the
first 60 seconds whether to engage or to look for holes.  The frame you set
in that first minute determines which mode they enter.  This section teaches
you to set the right frame.
\end{prereq}

There are two ways to open a quant presentation.

\begin{itemize}[leftmargin=1.5em]
  \item \textbf{Good:} ``He--Krishnamurthy (2013) predicts that intermediary
    balance-sheet constraints price risk across asset classes.  We tested
    this with daily SecDB risk data and found that factor-concentration
    changes predict next-week drawdowns with $\IC = 0.05$.''
  \item \textbf{Bad:} ``We fed 47 risk features into a gradient-boosted
    tree and it predicted returns.''
\end{itemize}

The first framing invites engagement.  It says: \emph{there is a
well-grounded economic hypothesis, and we used modern statistical tools to
test it.}  The audience evaluates the hypothesis, then checks whether the
test is sound.

The second invites skepticism.  It says: \emph{we did a fishing expedition
and found a pattern.}  The audience immediately asks ``how many features did
you try?'' and reaches for the Deflated Sharpe Ratio
(Chapter~\ref{ch:validation}).

\begin{keyidea}[The ML Is the Test, Not the Story]
Frame your work as testing a well-grounded economic hypothesis with modern
statistical tools.  The theory
(Chapter~\ref{ch:intermediary}) provides the \emph{why}; the ML
(Chapters~\ref{ch:regularized} and~\ref{ch:trees}) provides the
\emph{how precisely}.  Lead with the theory.  Save the model details for
the methodology slide.
\end{keyidea}

Kelly and Xiu (2023) survey the financial ML literature and note that the
most impactful papers---Gu, Kelly, and Xiu (2020), Freyberger, Neuhierl, and
Weber (2020), Feng, Giglio, and Xiu (2020)---all begin with an asset-pricing
question, not a modeling trick.  The ML is evaluated by whether it answers
the economic question better than existing methods, not by its raw $R^2$.

\paragraph{Practical framing template.}
For each signal family you present, use this three-sentence structure:

\begin{enumerate}[leftmargin=1.5em]
  \item \textbf{Theory:} ``[Paper] predicts that [mechanism] should cause
    [observable effect].''
  \item \textbf{Data advantage:} ``SecDB gives us [specific data] that
    external researchers cannot observe.''
  \item \textbf{Result:} ``We find [IC / Sharpe / direction], consistent
    [or inconsistent] with the hypothesis.''
\end{enumerate}

\begin{warning}[Never Lead with Model Complexity]
If the first thing your audience hears is ``we used a 500-tree LightGBM
ensemble with Bayesian hyperparameter optimization,'' you have lost them.
They will spend the next 19 minutes looking for overfitting rather than
listening to your results.  Complexity is a \emph{liability} until you have
established that simpler methods cannot do the job---which is exactly why you
always show the ridge baseline (Section~\ref{sec:presenting:ridge}).
\end{warning}

% ══════════════════════════════════════════════════════════════
\section{The Primary Metrics Triad}
\label{sec:presenting:triad}
% ══════════════════════════════════════════════════════════════

\begin{prereq}[What These Metrics Measure]
If you have not seen IC, Sharpe, or turnover before, review
Chapter~\ref{ch:backtesting} for formal definitions and computation.  This
section focuses on how to \emph{report} them, not how to compute them.
\end{prereq}

Every signal should be judged on three dimensions simultaneously.
Cherry-picking one is misleading.

\begin{definition}[The Metrics Triad]
\begin{enumerate}[leftmargin=1.5em]
  \item \textbf{Information Coefficient ($\IC$):} Does the signal predict
    direction?  Report the mean rank correlation between signal and
    subsequent return, plus its standard deviation across time.
    Example: $\IC = 0.04 \pm 0.12$ (monthly, cross-sectional).
  \item \textbf{Sharpe Ratio ($\SR$):} Does the signal make money?
    Report both the raw annualized Sharpe and the $\DSR$-adjusted value
    (Chapter~\ref{ch:validation}).
    Example: $\SR = 1.2$, $\DSR = 0.8$ after 35 trials.
  \item \textbf{Turnover:} Is it tradeable?  Report average daily
    one-way turnover and the breakeven transaction cost.
    Example: turnover = 18\%/day, breakeven cost = 12 bps.
\end{enumerate}
\end{definition}

\paragraph{Why all three matter.}  A signal with $\IC = 0.08$ but
daily turnover of 60\% is untradeable---transaction costs will eat the
alpha.  A signal with $\SR = 2.0$ but $\DSR = 0.3$ is probably overfit.
A signal with low turnover and a decent Sharpe but $\IC = 0.00$ is
just a static factor exposure, not a timing signal.

\begin{workedexample}{Reporting the Triad for VaR Utilization}
Suppose your VaR-utilization signal (Chapter~\ref{ch:features}) produces
the following backtest results after purged cross-validation
(Chapter~\ref{ch:validation}):

\begin{itemize}[leftmargin=1.5em]
  \item Mean $\IC$: $0.045$ (rank correlation of signal with next-5-day
    returns, averaged across 48 monthly cross-sections)
  \item Std of $\IC$: $0.11$ (so $\IC$-$t = 0.045 / (0.11 / \sqrt{48})
    \approx 2.83$)
  \item Raw $\SR$: $1.35$ (annualized, long-short quintile portfolio)
  \item $\DSR$: $0.91$ (after accounting for 28 experiments logged in
    your experiment tracker, Chapter~\ref{ch:validation})
  \item Daily one-way turnover: $14\%$
  \item Breakeven cost: $\text{Annual alpha} / (\text{turnover} \times
    252 \times 2) = 8.1\% / (0.14 \times 252 \times 2) \approx 11.5$~bps
\end{itemize}

You would report this as:

\medskip
\begin{center}
\begin{tabular}{lcc}
  \toprule
  \textbf{Metric} & \textbf{VaR Utilization} & \textbf{Ridge Baseline} \\
  \midrule
  $\IC$ (mean $\pm$ std) & $0.045 \pm 0.11$ & $0.031 \pm 0.13$ \\
  $\IC$-$t$              & 2.83              & 1.69 \\
  $\SR$ (annualized)     & 1.35              & 0.92 \\
  $\DSR$ (28 trials)     & 0.91              & 0.58 \\
  Turnover (\%/day)      & 14\%              & 9\% \\
  Breakeven cost (bps)   & 11.5              & 14.8 \\
  \bottomrule
\end{tabular}
\end{center}

\medskip
Key observations: (1)~the GBM beats ridge on IC and Sharpe, so the
nonlinearity adds value; (2)~the DSR remains positive even after 28
trials; (3)~turnover is moderate and the breakeven cost of 11.5~bps is
realistic for institutional execution.  The ridge baseline's higher
breakeven cost (14.8~bps) reflects lower turnover despite lower alpha---it
trades less but earns less.
\end{workedexample}

\begin{intuition}[The Triad as a Sanity Triangle]
Think of IC, Sharpe, and turnover as the vertices of a triangle.
A signal that scores well on all three is tradeable, predictive, and
profitable.  If any vertex collapses, the triangle breaks:
high IC + high Sharpe + high turnover = costs kill you;
high IC + low turnover + low Sharpe = the signal is real but the
effect size is too small to monetize;
high Sharpe + low IC = you are harvesting a risk premium, not predicting.
Always check all three.
\end{intuition}

% ══════════════════════════════════════════════════════════════
\section{Ridge Baseline on Every Chart}
\label{sec:presenting:ridge}
% ══════════════════════════════════════════════════════════════

\begin{prereq}[Why Ridge Is the Right Baseline]
Ridge regression (Chapter~\ref{ch:regularized}) is the simplest model
that can handle correlated features without blowing up.  It is linear,
fully interpretable, and hard to overfit.  If a gradient-boosted tree
cannot beat ridge on identical features, the nonlinearity is not helping
---the problem is either linear or your tree is overfitting.  This was
discussed in detail in Chapter~\ref{ch:regularized},
Section~\ref{sec:regularized-summary}.
\end{prereq}

\paragraph{The rule.}  Every cumulative P\&L chart, every IC time series,
every Sharpe bar chart has \emph{two} lines:

\begin{itemize}[leftmargin=1.5em]
  \item \textbf{Solid line:} your model (LightGBM or whatever you used)
  \item \textbf{Dashed line:} ridge baseline on the identical feature set
\end{itemize}

Three possible outcomes:

\begin{enumerate}[leftmargin=1.5em]
  \item \textbf{Clear separation (GBM $\gg$ ridge).}  The nonlinearity
    adds value.  Show where the lines diverge and use SHAP
    (Section~\ref{sec:presenting:shap}) to explain \emph{why} the tree
    found something ridge missed---typically an interaction or a threshold
    effect.
  \item \textbf{Lines overlap (GBM $\approx$ ridge).}  The problem is
    approximately linear.  This is \emph{not} a failure---it means your
    features are good and the relationship is simple.  Say so honestly:
    ``Ridge captures the signal; the GBM adds marginal improvement that
    does not survive transaction costs.''  Then use ridge as your
    production model.  It is simpler, faster, and more robust.
  \item \textbf{Ridge wins (GBM $<$ ridge).}  The tree overfit.  Do
    not present the tree results.  Present ridge, explain why you tried
    trees, and report that they did not help.  This is an honest negative
    result (Section~\ref{sec:presenting:negatives}).
\end{enumerate}

\begin{keyidea}[The Ridge Baseline Is Your Credibility Insurance]
Showing the ridge baseline on every chart is the single most effective
way to build credibility with a quantitative audience.  It signals three
things: (1)~you understand the bias-variance tradeoff; (2)~you are not
hiding behind model complexity; (3)~you would have reported honestly if
the simple model won.  Kelly and Xiu (2023) note that the most
convincing ML-in-finance papers are those that clearly demonstrate when
and why ML adds value over linear benchmarks.
\end{keyidea}

% ══════════════════════════════════════════════════════════════
\section{SHAP Waterfall Plots}
\label{sec:presenting:shap}
% ══════════════════════════════════════════════════════════════

\begin{prereq}[SHAP Background]
Chapter~\ref{ch:interpretation} covers SHAP values in detail:
Shapley's axioms, TreeSHAP computation, and the difference between
global importance (mean $|\SHAP|$) and local explanation (waterfall).
This section focuses specifically on how to \emph{present} waterfall
plots to a non-technical audience.
\end{prereq}

For the top 3 most interesting predictions your model makes, show a SHAP
waterfall plot.  These convert ``trust me, the model works'' into ``here
is exactly what drove this prediction.''

\paragraph{How to read a waterfall plot.}  A waterfall plot for a single
observation shows:

\begin{enumerate}[leftmargin=1.5em]
  \item The \textbf{base value} $\E[f(\bX)]$: the model's average
    prediction across the dataset.
  \item Each feature's \textbf{contribution}: a bar pushing the prediction
    up (red) or down (blue) from the base value.
  \item The \textbf{final prediction} $f(\bx_i)$: the sum of the base
    value and all feature contributions.
\end{enumerate}

\paragraph{Presenting waterfalls to a desk audience.}
Pick three predictions that illustrate different mechanisms:

\begin{enumerate}[leftmargin=1.5em]
  \item \textbf{A strong correct prediction.}  ``On 2024-03-15, the model
    predicted a 2.1\% drawdown in IG credit.  The main driver was a spike
    in factor-concentration ($\HHI$ up 40\% week-over-week,
    Chapter~\ref{ch:features}).  The drawdown materialized at 1.8\%.''
  \item \textbf{A strong wrong prediction.}  ``On 2024-07-02, the model
    predicted a drawdown that did not happen.  SHAP shows the driver was
    VaR utilization at 93\%---but a Fed announcement reversed the move.
    The model cannot see policy announcements.''
  \item \textbf{A prediction where GBM and ridge disagree.}  ``Ridge
    predicted no move; the GBM predicted a reversal.  SHAP shows the
    driver was an interaction between factor concentration and cross-asset
    flow (Chapter~\ref{ch:features})---a nonlinearity that ridge cannot
    capture.''
\end{enumerate}

\begin{warning}[Do Not Show 20 Waterfall Plots]
One well-chosen waterfall is worth more than twenty.  If you show too many,
the audience loses the thread.  Pick three.  Have backup slides with
additional examples in case someone asks ``what about [date]?''  The
global summary plot (mean $|\SHAP|$ bar chart) belongs in the methodology
section, not the results section.
\end{warning}

\paragraph{Explaining to non-technical audiences.}
Avoid the word ``Shapley.''  Say: ``This chart breaks down the model's
prediction into the contribution of each input.  Red bars push the
prediction up; blue bars push it down.  The biggest bar tells you what
mattered most for this particular day.''  That is sufficient.  If someone
asks how SHAP is computed, you can explain the game-theory foundation
(Chapter~\ref{ch:interpretation}), but do not volunteer it.

% ══════════════════════════════════════════════════════════════
\section{One Slide Per Claim}
\label{sec:presenting:slides}
% ══════════════════════════════════════════════════════════════

\begin{prereq}[Slide Design Principles]
The ``one slide per claim'' principle comes from scientific presentation
best practices (Tufte, 2001; Schwabish, 2014).  Each slide should
communicate exactly one takeaway.  If the audience remembers only the
title of each slide, they should understand your full argument.
\end{prereq}

\paragraph{Slide structure.}
Every slide follows the same template:

\begin{itemize}[leftmargin=1.5em]
  \item \textbf{Title = claim.}  Not ``Results'' but ``Factor
    concentration predicts IG drawdowns with IC = 0.05.''  The title is
    the takeaway.
  \item \textbf{Body = one chart.}  The chart proves the claim.  No
    multi-chart slides---they split attention and invite the audience to
    look at the wrong chart while you are talking about the right one.
  \item \textbf{Footer = metric.}  A single line with the key number:
    ``$\IC = 0.045$, $\SR = 1.35$, $\DSR = 0.91$, $n = 1{,}247$.''
\end{itemize}

\paragraph{Suggested 10-slide deck.}
For the Phase~5 presentation (Task~24), structure your deck as follows:

\begin{enumerate}[leftmargin=1.5em]
  \item \textbf{Title + thesis.}  ``Risk as Alpha: Testing Intermediary
    Asset Pricing with Daily SecDB Data.''  One sentence summarizing the
    core finding.
  \item \textbf{Theory.}  He--Krishnamurthy (2013), Adrian--Etula--Muir
    (2014): dealer constraints price risk.  SecDB gives daily data where
    academics use quarterly.  (Chapter~\ref{ch:intermediary})
  \item \textbf{Data and methodology.}  Feature families
    (Chapter~\ref{ch:features}), validation framework
    (Chapter~\ref{ch:validation}), backtesting with costs
    (Chapter~\ref{ch:backtesting}).  One flow diagram.
  \item \textbf{Result: Signal Family A.}  Cumulative P\&L chart with
    ridge baseline.  Footer: IC, Sharpe, DSR, turnover.
  \item \textbf{Result: Signal Family B.}  Same format.
  \item \textbf{Result: Signal Family C.}  Same format.  (If only two
    families survived, use this slide for the combined model.)
  \item \textbf{Ridge vs.\ GBM comparison.}  Side-by-side bar chart of
    IC and Sharpe for each signal family.  Footer: ``GBM adds [X]\% IC
    over ridge.''
  \item \textbf{Regime decomposition.}  Signal Sharpe by regime
    (Chapter~\ref{ch:regimes}).  Footer: ``Signal works in [N] of [K]
    regimes.''
  \item \textbf{Capacity and costs.}  Sharpe vs.\ transaction cost curve
    (Chapter~\ref{ch:backtesting}).  Footer: breakeven cost.
  \item \textbf{Next steps.}  What you would do with 6 more months.
    Extensions, additional data, live paper trading.
\end{enumerate}

\begin{keyidea}[If They Only Read the Titles]
Write your slide titles so that reading just the titles in sequence tells
the complete story.  Test this: list your 10 titles.  If a colleague can
understand your argument from the titles alone, the deck is well
structured.  If they cannot, rewrite the titles until they can.
\end{keyidea}

% ══════════════════════════════════════════════════════════════
\section{What Didn't Work as Credibility}
\label{sec:presenting:negatives}
% ══════════════════════════════════════════════════════════════

Documented negative results are your strongest credibility signal.  A
presentation that reports only successes triggers the audience's
overfitting alarm.  One that says ``we tested five signal families;
three worked, two didn't'' signals intellectual honesty.

\paragraph{What counts as a valuable negative result.}

\begin{itemize}[leftmargin=1.5em]
  \item ``Cross-asset flow features (Chapter~\ref{ch:features}) showed
    no predictive power for VIX innovations after purged CV
    (Chapter~\ref{ch:validation}).  $\IC = 0.003 \pm 0.14$, not
    significantly different from zero.''
  \item ``Scenario P\&L dispersion predicted drawdowns in-sample
    ($\IC = 0.06$) but failed walk-forward validation
    (Chapter~\ref{ch:backtesting}).  We attribute this to a structural
    break in the scenario set after 2022.''
  \item ``The GBM did not beat ridge on VaR-dynamics features.  The
    relationship appears to be linear.''
\end{itemize}

Each of these is genuinely informative.  They tell the audience what the
\emph{data} look like, not just what the \emph{model} found.

\paragraph{Where to put negatives in your presentation.}
Include one slide between the results slides and the regime decomposition
slide (after slide 6, before slide 8 in the suggested deck).  Title it
honestly: ``Signal families that did not survive validation.''  Show a
table:

\medskip
\begin{center}
\begin{tabular}{lccl}
  \toprule
  \textbf{Feature Family} & \textbf{IC (purged CV)} & \textbf{DSR} &
    \textbf{Why It Failed} \\
  \midrule
  Cross-asset flow & $0.003$ & $-0.12$ & No signal after detrending \\
  Scenario P\&L    & $0.061$ (IS) / $0.008$ (OOS) & $0.04$ &
    Structural break in scenarios \\
  \bottomrule
\end{tabular}
\end{center}

\begin{intuition}[Negatives Build Trust Because They Are Costly]
Reporting negatives is costly: it uses scarce presentation time and
shows the limits of your work.  The audience knows this.  The fact that
you are \emph{willing} to pay that cost signals that the positives are
real.  A presenter who hides failures is optimizing for impression, not
information.  A desk head can tell the difference.
\end{intuition}

% ══════════════════════════════════════════════════════════════
\section{Handling Desk Questions}
\label{sec:presenting:questions}
% ══════════════════════════════════════════════════════════════

\begin{prereq}[Know Your Backup Slides]
For every question below, have a backup slide with the detailed chart or
table.  Your verbal answer should be two sentences.  If they want more,
flip to the backup slide.  Prepare 5--8 backup slides covering the
scenarios below.
\end{prereq}

The following questions come up in almost every desk presentation of a
quantitative signal.  Prepare two-sentence answers for each.

\paragraph{``What's the capacity?''}
\emph{Answer:} ``At current turnover of [X]\%/day, the signal breaks
even at [Y]~bps per side.  Institutional execution costs for [asset
class] are approximately [Z]~bps, so estimated capacity is
\$[N]~million before market impact becomes material.''

\noindent\emph{Backup slide:} Sharpe ratio vs.\ transaction cost curve
(Chapter~\ref{ch:backtesting}).  Show the point where Sharpe hits zero.

\paragraph{``What happens in a crisis?''}
\emph{Answer:} ``We decomposed performance by regime
(Chapter~\ref{ch:regimes}).  The signal earned $\SR = [A]$ in
steady-state periods and $\SR = [B]$ during stress.  It [does / does
not] work in both regimes.''

\noindent\emph{Backup slide:} Regime-conditional cumulative P\&L chart
with GMM regime shading.

\paragraph{``Why not just a linear model?''}
\emph{Answer:} ``We tested ridge on identical features.  The GBM
improves IC by [X]\% and Sharpe by [Y]\%.  SHAP identifies two
interaction effects that ridge cannot capture: [feature A $\times$
feature B] and [threshold in feature C].''

\noindent\emph{Backup slide:} Ridge vs.\ GBM bar chart
(Section~\ref{sec:presenting:ridge}) plus SHAP dependence plot for the
key interaction (Chapter~\ref{ch:interpretation}).

\paragraph{``How is this different from [public factor]?''}
\emph{Answer:} ``We ran confound checks by regressing our signal on
[Fama--French / Barra / known macro factors].  The residual IC is
[X], which means [Y]\% of the signal is orthogonal to published
factors.''

\noindent\emph{Backup slide:} Table of alpha after controlling for known
factors (Chapter~\ref{ch:panel}).

\paragraph{``Is this overfit?''}
\emph{Answer:} ``We tracked [N] experiments in total.  The Deflated
Sharpe Ratio is [X], which means the result survives correction for
multiple testing.  Walk-forward out-of-sample Sharpe is [Y], consistent
with purged CV estimates.''

\noindent\emph{Backup slide:} CPCV distribution plot
(Chapter~\ref{ch:validation}) and walk-forward equity curve.

\begin{warning}[Never Say ``I Don't Know'' Without a Follow-Up]
If you genuinely do not know the answer, say: ``I haven't tested that
specific scenario.  My best estimate based on [related analysis] is [X].
I can run that test and follow up by [date].''  This is honest, shows
you understand what test would answer the question, and commits to a
deliverable.  Never just say ``I don't know'' and move on.
\end{warning}

% ══════════════════════════════════════════════════════════════
\section{Signal Decay and Durability}
\label{sec:presenting:decay}
% ══════════════════════════════════════════════════════════════

\begin{prereq}[Why Signals Decay]
Signal decay is the phenomenon where a predictive pattern weakens over
time.  This happens for two reasons: (1)~\emph{statistical decay}---the
pattern was partially spurious (data mining); (2)~\emph{economic
decay}---other market participants discover and trade the pattern,
arbitraging it away.  McLean and Pontiff (2016) disentangle these two
channels.
\end{prereq}

\begin{keyresult}[McLean--Pontiff: Published Signals Lose 26\% OOS,
  58\% Post-Publication]
McLean and Pontiff (2016, \textit{Journal of Finance}) study 97
cross-sectional return predictors documented in academic publications.
They find:
\begin{itemize}[leftmargin=1.5em]
  \item \textbf{Out-of-sample decay:} Portfolio returns are \textbf{26\%
    lower} in the post-sample period (after the original paper's sample
    ends but before publication).  This is an upper bound on the
    data-mining component.
  \item \textbf{Post-publication decay:} Returns decline by a further
    \textbf{32\%} (58\% total) after the paper is published.
    This additional decay ($58\% - 26\% = 32\%$) is attributed to
    publication-informed trading: arbitrageurs read the paper and trade
    the signal.
  \item \textbf{Implication:} Roughly half of the apparent decay is real
    (arbitrage), and half is data mining.
\end{itemize}
\end{keyresult}

\paragraph{What this means for your project.}
Your project uses proprietary SecDB data.  This creates a partial shield
against the McLean--Pontiff decay mechanism:

\begin{enumerate}[leftmargin=1.5em]
  \item \textbf{Competitors cannot replicate internal risk signals.}
    External researchers approximate dealer positions from public data
    (13F filings, CFTC reports, Fed Z.1).  These proxies are quarterly,
    lagged, and noisy (Chapter~\ref{ch:intermediary}).  SecDB data is
    daily and precise.  If your signal's alpha comes from the
    \emph{timeliness} of the data, competitors with quarterly data
    cannot replicate it.
  \item \textbf{But correlation with published factors reduces the
    shield.}  If your VaR-utilization signal is 80\% correlated with
    Adrian--Etula--Muir's (2014) intermediary leverage factor, then 80\%
    of the signal shares that factor's post-publication decay.  The
    confound checks (Chapter~\ref{ch:panel}) directly measure this
    overlap.
  \item \textbf{Structural signals decay slower.}  A signal based on
    regulatory constraints (VaR limits, capital requirements) decays
    only if the regulation changes.  Basel III has been stable since 2013.
    A signal based on behavioral patterns (crowding, herding) decays
    faster because behavior adapts.
\end{enumerate}

\begin{keyidea}[Be Honest About Durability]
A signal that decays is still valuable if you are the first to trade it.
The question is not ``will this signal last forever?'' but ``how long is
the edge, and can we monetize it before it decays?''  Present a realistic
assessment: ``Based on the proprietary data advantage and low correlation
with published factors ($r = [X]$), we estimate a useful life of [Y]
years before competitors can approximate the signal with public data.''
\end{keyidea}

\paragraph{Presenting decay honestly.}
In your report (Task~23) and presentation (Task~24), include a brief
discussion of durability.  Frame it around three questions:

\begin{enumerate}[leftmargin=1.5em]
  \item \textbf{Data advantage:} Can an external researcher replicate
    this signal from public data?  If yes, the McLean--Pontiff 58\%
    decay applies in full.
  \item \textbf{Factor overlap:} What fraction of the signal's variance
    is explained by published factors?  Report the $R^2$ from the
    confound regression (Chapter~\ref{ch:panel}).
  \item \textbf{Structural vs.\ behavioral:} Is the signal driven by
    regulatory constraints (slow to change) or trader behavior (fast to
    adapt)?
\end{enumerate}

% ══════════════════════════════════════════════════════════════
\section{The Research Report}
\label{sec:presenting:report}
% ══════════════════════════════════════════════════════════════

Task~23 calls for a full desk-ready research report.  This is a written
document, not slides.  Structure it as follows:

\begin{enumerate}[leftmargin=1.5em]
  \item \textbf{Executive summary} (1 page).  Thesis, data, key result,
    capacity estimate.  A busy desk head reads only this page.
  \item \textbf{Hypothesis and motivation} (1--2 pages).  Economic
    theory (Chapter~\ref{ch:intermediary}), data advantage, prior
    literature.
  \item \textbf{Data and features} (2--3 pages).  Feature families
    (Chapter~\ref{ch:features}), point-in-time construction
    (Chapter~\ref{ch:backtesting}), sample period, holdout reservation.
  \item \textbf{Methodology} (2--3 pages).  Labeling
    (Chapter~\ref{ch:labeling}), validation framework
    (Chapter~\ref{ch:validation}), backtesting engine
    (Chapter~\ref{ch:backtesting}), model specifications.
  \item \textbf{Results} (3--5 pages).  One subsection per signal family.
    Each subsection: IC table, cumulative P\&L chart (with ridge
    baseline), SHAP summary, regime decomposition.
  \item \textbf{Negative results} (1 page).  What did not work and why.
  \item \textbf{Robustness} (1--2 pages).  CPCV distribution, DSR,
    walk-forward OOS, rolling IC stability, confound checks.
  \item \textbf{Capacity and costs} (1 page).  Sharpe vs.\ cost curve,
    breakeven cost, estimated dollar capacity.
  \item \textbf{Limitations and next steps} (1 page).  Sample size
    caveats, data coverage gaps, suggested extensions.
\end{enumerate}

\paragraph{One chart per claim.}  Every chart in the report should
support exactly one claim, stated in the chart's caption.  Bad caption:
``Figure 3: Results.''  Good caption: ``Figure 3: Factor-concentration
HHI predicts IG credit drawdowns ($\IC = 0.05$, $\DSR = 0.91$).  Solid
line: LightGBM; dashed line: ridge baseline.''

\begin{warning}[Do Not Iterate After Unblinding the Holdout]
Chapter~\ref{ch:backtesting} made this point, but it bears repeating in
the presentation context.  The walk-forward OOS test on the reserved
holdout (Task~22) is a \emph{one-shot} exercise.  You run it once,
report the result, and do not go back to tune the model.  If the OOS
result is disappointing, report it honestly.  ``In-sample $\SR = 1.35$,
out-of-sample $\SR = 0.72$'' is a legitimate result.  ``We re-optimized
on the holdout until OOS Sharpe improved'' is fraud.
\end{warning}

% ══════════════════════════════════════════════════════════════
\section{Common Presentation Mistakes}
\label{sec:presenting:mistakes}
% ══════════════════════════════════════════════════════════════

\begin{enumerate}[leftmargin=1.5em]
  \item \textbf{Leading with $R^2$.}  In finance, $R^2$ values are
    tiny by design.  Gu, Kelly, and Xiu (2020) report monthly
    out-of-sample $R^2$ of 0.40\% for their best model.  If you lead
    with ``$R^2 = 0.003$,'' your audience hears ``this model explains
    nothing.''  Lead with IC and Sharpe instead.
  \item \textbf{Reporting in-sample metrics only.}  Every metric you
    present should be out-of-sample (purged CV or walk-forward).
    In-sample metrics belong in the appendix, if anywhere.
  \item \textbf{Omitting the number of trials.}  If you report
    $\SR = 1.2$ without saying how many experiments you ran, the
    audience cannot assess whether the result survives multiple-testing
    correction.  Bailey, Borwein, and Lopez de Prado (2014) show that
    even with a true $\SR = 0$, running $\sim$45 independent trials on
    5 years of daily data will produce a maximum observed $\SR$ near 1.0
    by pure chance.  Always report the trial count and the $\DSR$.
  \item \textbf{Showing only the best signal.}  If you tested five
    signal families and show only the one that worked, you are
    p-hacking by omission.  Show all five.  The ones that failed are
    credibility (Section~\ref{sec:presenting:negatives}).
  \item \textbf{Ignoring transaction costs.}  A signal with $\SR = 2.0$
    gross and $\SR = 0.3$ net is not a good signal.  Always report net
    Sharpe.  Chapter~\ref{ch:backtesting} covers cost modeling in
    detail.
  \item \textbf{Confusing statistical significance with economic
    significance.}  An $\IC$-$t$ of 3.5 on a signal that generates
    20~bps/year of alpha is statistically significant but economically
    irrelevant for a desk that needs 200~bps to justify the
    infrastructure.
\end{enumerate}

% ══════════════════════════════════════════════════════════════
\section{Tying It All Together}
\label{sec:presenting:synthesis}
% ══════════════════════════════════════════════════════════════

This learning guide has covered a complete research pipeline:

\begin{enumerate}[leftmargin=1.5em]
  \item \textbf{Theory} (Chapters~\ref{ch:asset-pricing}
    and~\ref{ch:intermediary}):  Why dealer balance-sheet constraints
    should predict returns.
  \item \textbf{Market context}
    (Chapters~\ref{ch:risk-systems} and~\ref{ch:microstructure}):
    How risk systems and dealer positioning create the data you observe.
  \item \textbf{Statistical methods}
    (Chapters~\ref{ch:regularized}, \ref{ch:trees},
    and~\ref{ch:panel}):  Ridge, LightGBM, and panel regression---the
    tools you use to test the hypothesis.
  \item \textbf{Interpretation} (Chapter~\ref{ch:interpretation}):
    SHAP, MDI, MDA---how to explain \emph{why} the model predicts what
    it predicts.
  \item \textbf{Feature engineering} (Chapter~\ref{ch:features}):
    Turning raw SecDB outputs into testable signals.
  \item \textbf{Labeling} (Chapter~\ref{ch:labeling}):  Triple-barrier,
    meta-labeling, and standard return targets.
  \item \textbf{Validation} (Chapter~\ref{ch:validation}):  Purged CV,
    CPCV, DSR, Haircut Sharpe---the honest scorecard.
  \item \textbf{Backtesting} (Chapter~\ref{ch:backtesting}):
    Transaction-cost-aware P\&L and walk-forward OOS.
  \item \textbf{Regimes} (Chapter~\ref{ch:regimes}):  Is your signal
    robust across market conditions?
  \item \textbf{Presentation} (this chapter):  How to communicate
    all of the above so that the work speaks for itself.
\end{enumerate}

The connecting thread is the same thesis that appeared in Chapter~\ref{ch:intermediary}:
intermediary constraints price risk, SecDB data captures those constraints
daily, and the question is whether the data advantage translates into
tradeable alpha.  Everything else---the infrastructure, the modeling, the
validation, the presentation---is in service of answering that question
rigorously.

% ══════════════════════════════════════════════════════════════
\section{Summary}
\label{sec:presenting:summary}
% ══════════════════════════════════════════════════════════════

\begin{center}
\begin{tabular}{p{4cm} p{9.5cm}}
  \toprule
  \textbf{Concept} & \textbf{Key Takeaway} \\
  \midrule
  Framing &
    Lead with the economic hypothesis, not the ML model.  The ML is the
    test, not the story. \\[4pt]
  Metrics triad &
    Report IC (predictiveness), Sharpe (profitability), and turnover
    (tradeability) together.  Cherry-picking one is misleading. \\[4pt]
  Ridge baseline &
    Every chart has two lines: model and ridge.  If ridge wins, say so.
    This is your credibility insurance. \\[4pt]
  SHAP waterfalls &
    Pick three predictions.  Show what drove each one.  Convert ``trust
    me'' into ``here's the mechanism.'' \\[4pt]
  One slide per claim &
    Title = claim, body = one chart, footer = metric.  10 slides total.
    If you read only the titles, the story should be clear. \\[4pt]
  Negative results &
    Report what didn't work.  It shows intellectual honesty and
    prevents the audience from wondering what you are hiding. \\[4pt]
  Desk Q\&A &
    Prepare two-sentence answers for: capacity, crisis behavior,
    linear vs.\ nonlinear, public factor overlap, overfitting.
    Have backup slides with detailed charts. \\[4pt]
  Signal decay &
    McLean--Pontiff: 26\% OOS decay, 58\% post-publication.
    Proprietary data partially shields against publication-informed
    trading, but factor overlap reduces the shield. \\[4pt]
  Research report &
    Structured document: executive summary, hypothesis, data,
    methodology, results (with ridge), negatives, robustness, capacity,
    next steps.  One chart per claim. \\[4pt]
  Common mistakes &
    Do not lead with $R^2$.  Report OOS metrics.  State the trial count.
    Show all signals tested.  Report net-of-cost Sharpe. \\
  \bottomrule
\end{tabular}
\end{center}
